%Daten zur Institution

%\input{Dozentdaten}

%\renewcommand{\fachbereich}{Fachbereich}

%\renewcommand{\dozent}{Prof. Dr. . }

%Klausurdaten

\renewcommand{\klausurgebiet}{ }

\renewcommand{\klausurtyp}{ }

\renewcommand{\klausurdatum}{ . 20}

\klausurvorspann {\fachbereich} {\klausurdatum} {\dozent} {\klausurgebiet} {\klausurtyp}

%Daten für folgende Punktetabelle

\renewcommand{\aeins}{  3  }

\renewcommand{\azwei}{  3   }

\renewcommand{\adrei}{  7  }

\renewcommand{\avier}{  4  }

\renewcommand{\afuenf}{  0  }

\renewcommand{\asechs}{  0  }

\renewcommand{\asieben}{  0  }

\renewcommand{\aacht}{  5  }

\renewcommand{\aneun}{  15  }

\renewcommand{\azehn}{  5  }

\renewcommand{\aelf}{  0  }

\renewcommand{\azwoelf}{  3  }

\renewcommand{\adreizehn}{  3   }

\renewcommand{\avierzehn}{  0  }

\renewcommand{\afuenfzehn}{  10  }

\renewcommand{\asechzehn}{  4  }

\renewcommand{\asiebzehn}{  62  }

\renewcommand{\aachtzehn}{    }

\renewcommand{\aneunzehn}{    }

\renewcommand{\azwanzig}{    }

\renewcommand{\aeinundzwanzig}{    }

\renewcommand{\azweiundzwanzig}{    }

\renewcommand{\adreiundzwanzig}{    }

\renewcommand{\avierundzwanzig}{    }

\renewcommand{\afuenfundzwanzig}{    }

\renewcommand{\asechsundzwanzig}{    }

\punktetabellesechzehn

\klausurnote

\newpage

\setcounter{section}{K}

__NOEDITSECTION__

\inputaufgabeklausurloesung
{3}

{

Definiere die folgenden
\zusatzklammer {kursiv gedruckten} {} {} Begriffe.
\aufzaehlungsechs{Eine
\stichwort {bogenparametrisierte} {}
Kurve
\maabbdisp {\gamma} {[a,b]} { \R^n
} {.}

}{Zwei in einem Punkt

\mavergleichskette
{\vergleichskette
{  P
}
{ \in }{ M
}
{  }{}
{    }{}
{   }{}
}
{}{}{}
einer differenzierbaren Mannigfaltigkeit $M$ \stichwort {tangential äquivalente} {} differenzierbare Kurven
\maabbdisp {\gamma_1, \gamma_2} { I } { M
} {}
\zusatzklammer {dabei sei

\mavergleichskettek
{\vergleichskettek
{  0
}
{ \in }{  I
}
{  }{
}
{    }{
}
{   }{
}
}
{}{}{}
ein offenes reelles Intervall und

\mavergleichskettek
{\vergleichskettek
{   \gamma_1(0)
}
{ = }{ \gamma_2(0)
}
{ = }{ P
}
{    }{
}
{   }{
}
}
{}{}{}} {} {.}

}{Ein
\stichwort {stetiger Schnitt} {}
zu einer stetigen Abbildung
\maabbdisp {p} {Y} {X
} {}
zwischen topologischen Räumen
\mathkor {} {X} {und} {Y} {.}

}{Die \stichwort {kanonische Volumenform} {} auf einer orientierten riemannschen Mannigfaltigkeit $M$.

}{Die
\stichwort {äußere Ableitung} {}
einer
\definitionsverweis {stetig differenzierbaren}{}{}
$k$-\definitionsverweis {Differentialform}{}{}

\mathl{\omega \in
{ \mathcal E }^{ k } (  U   )}{} auf einer
\definitionsverweis {offenen Menge}{}{}

\mathl{U \subseteq \R^n}{.}

}{Die \stichwort {Schnittkrümmung} {} zu linear unabhängigen Vektoren

\mavergleichskette
{\vergleichskette
{  v,w
}
{ \in }{T_PM
}
{  }{
}
{    }{
}
{   }{
}
}
{}{}{}
auf einer
\definitionsverweis {riemannschen Mannigfaltigkeit}{}{}
$M$ in einem Punkt

\mavergleichskette
{\vergleichskette
{ P
}
{ \in }{ M
}
{  }{
}
{    }{
}
{   }{
}
}
{}{}{.}
}

}
{

\aufzaehlungsechs{Eine
\definitionsverweis {differenzierbare Kurve}{}{}
\maabbdisp {\gamma} {[a,b]} {\R^n
} {,}
heißt
bogenparametrisiert,
wenn

\mavergleichskette
{\vergleichskette
{ \gamma_1'(t)^2 + \cdots + \gamma_n'(t)^2
}
{ = }{ 1
}
{  }{
}
{    }{
}
{   }{
}
}
{}{}{}
für alle $t$ gilt.
}{Die beiden Kurven
\mathkor {} {\gamma_1} {und} {\gamma_2} {}
heißen tangential äquivalent in $P$, wenn es eine offene Umgebung
\mathl{P \in U}{} und eine
\definitionsverweis {Karte}{}{}
\maabbdisp {\alpha} { U } { V
} {}
mit
\mathl{V \subseteq \R^n}{} derart gibt, dass

\mavergleichskettedisp
{\vergleichskette
{  \left(  \alpha \circ \left(  \gamma_1 \vert_{\gamma_1^{-1} (U)}  \right)   \right)'(0)
}
{ =} { \left(  \alpha \circ \left(  \gamma_2 \vert_{\gamma_2^{-1} (U)}  \right)  \right)'(0)
}
{ } {
}
{ } {
}
{ } {
}
}
{}{}{} gilt.
}{Unter einem
stetigen Schnitt
zu $p$ versteht man eine stetige Abbildung
\maabb {s} {X} {Y
} {}
mit

\mavergleichskettedisp
{\vergleichskette
{  p \circ s
}
{ =} {
\operatorname{Id}_{ X }
}
{ } {
}
{ } {
}
{ } {
}
}
{}{}{.}
}{Zu
\mathl{P \in M}{} sei
\mathl{\omega_P}{} diejenige
\definitionsverweis {alternierende Form}{}{}
auf
\mathl{T_PM}{}
\zusatzklammer {bzw. das entsprechende Element aus \mathlk{\bigwedge^n T^*_PM}{}} {} {,}
die jeder die
\definitionsverweis {Orientierung}{}{}
repräsentierenden
\definitionsverweis {Orthonormalbasis}{}{}
den Wert $1$ zuordnet. Dann heißt die
$n$-\definitionsverweis {Differentialform
}{}{}
\maabbeledisp {} { M } { \bigwedge^n T^*M
} {P } { \omega_P
} {,}
die kanonische Volumenform auf $M$.
}{Die Form $\omega$ besitzt auf $U$ eine Darstellung

\mathdisp {\omega = \sum_{I \subseteq \{1 , \ldots , n \},\,  { \# \left(   I  \right) } =k } f_I dx_I} {  }

mit
\definitionsverweis {stetig differenzierbaren Funktionen}{}{}
\maabbdisp {f_I} { U } {\R
} {.}
Dann ist die äußere Ableitung die $(k+1)$-Form

\mavergleichskettedisp
{\vergleichskette
{ d \omega
}
{ =} { \sum_{I \subseteq \{1 , \ldots , n \},\,  { \# \left(   I  \right) } = k } df_I \wedge dx_I
}
{ =} { \sum_{I \subseteq \{1 , \ldots , n \},\,  { \# \left(   I  \right) } = k } \left(  \sum_{j = 1}^n { \frac{   \partial  f_I  }{  \partial  x_j  } } dx_j  \right)  \wedge dx_I
}
{ } {
}
{ } {
}
}
{}{}{.}
}{Man nennt die Zahl

\mavergleichskettedisp
{\vergleichskette
{ K(v,w)
}
{ =} {  { \frac{    \left\langle R(v,w)w , v  \right\rangle  }{   \left\langle v , v  \right\rangle  \left\langle w , w  \right\rangle  -\left\langle v , w  \right\rangle^2   } }
}
{ } {
}
{ } {
}
{ } {
}
}
{}{}{,}
wobei $R$ den
\definitionsverweis {Krümmungsoperator}{}{}
auf dem Tangentialbündel bezeichnet, die
\definitionswort {Schnittkrümmung}{}
zu $v,w$ in $P$.
}

 }

__NOEDITSECTION__

\inputaufgabeklausurloesung
{3}

{

Formuliere die folgenden Sätze.
\aufzaehlungdrei{Der
\stichwort {Satz über die Orientierungen auf einer differenzierbaren Hyperfläche} {}

\mavergleichskette
{\vergleichskette
{  Y
}
{ \subseteq }{  \R^n
}
{  }{
}
{    }{
}
{   }{
}
}
{}{}{.}}{Die Formel für die Lie-Klammer von Vektorfeldern auf einer offenen Menge

\mavergleichskette
{\vergleichskette
{  U
}
{ \subseteq }{  \R^n
}
{  }{
}
{    }{
}
{   }{
}
}
{}{}{.}}{Der Satz über
\stichwort {Retraktionen zum Rand} {}
auf Mannigfaltigkeiten.}

}
{

\aufzaehlungdrei{\faktsituation {Es sei

\mavergleichskette
{\vergleichskette
{ W
}
{ \subseteq }{ \R^n
}
{  }{
}
{    }{
}
{   }{
}
}
{}{}{}
\definitionsverweis {offen}{}{,}
\maabb {h} { W } { \R
} {}
eine
\definitionsverweis {stetig differenzierbare Funktion}{}{}
und

\mavergleichskette
{\vergleichskette
{ Y
}
{ = }{ h^{-1}(c)
}
{  }{
}
{    }{
}
{   }{
}
}
{}{}{}
die
\definitionsverweis {Faser}{}{}
zu

\mavergleichskette
{\vergleichskette
{ c
}
{ \in }{ \R
}
{  }{
}
{    }{
}
{   }{
}
}
{}{}{,}
wobei $h$ in jedem Punkt von $Y$
\definitionsverweis {regulär}{}{}
sei. Es sei $Y$
\definitionsverweis {zusammenhängend}{}{.}}
\faktfolgerung {Dann gibt es genau zwei
\definitionsverweis {Orientierungen}{}{}
auf $Y$.}
\faktzusatz {}
\faktzusatz {}}{\faktsituation {Es sei

\mavergleichskette
{\vergleichskette
{ U
}
{ \subseteq }{ \R^n
}
{  }{
}
{    }{
}
{   }{
}
}
{}{}{}
\definitionsverweis {offen}{}{}
und seien $f_1 , \ldots , f_n, g_1 , \ldots , g_n$ zweifach
\definitionsverweis {stetig differenzierbare Funktionen}{}{}
auf $U$ mit den zugehörigen Differentialoperatoren

\mavergleichskette
{\vergleichskette
{ D
}
{ = }{ \sum_{i = 1}^n f_i \partial_i
}
{  }{
}
{    }{
}
{   }{
}
}
{}{}{}
und

\mavergleichskette
{\vergleichskette
{ E
}
{ = }{ \sum_{j = 1}^n g_j \partial_j
}
{  }{
}
{    }{
}
{   }{
}
}
{}{}{}}
\faktfolgerung {Dann gilt

\mavergleichskettedisp
{\vergleichskette
{ D \circ E -E \circ D
}
{ =} { \sum_{j = 1}^n \left(  \sum_{i = 1}^n \left(  f_i  \partial_i g_j - g_i \partial_i f_j  \right)  \right)  \partial_j
}
{ } {
}
{ } {
}
{ } {
}
}
{}{}{.}
Insbesondere entspricht dieser Ausdruck selbst einem Differentialoperator der Ordnung $1$ und somit einem Vektorfeld.}
\faktzusatz {}
\faktzusatz {}}{Es sei $M$ eine
\definitionsverweis {kompakte}{}{}
\definitionsverweis {orientierte}{}{}
\definitionsverweis {differenzierbare Mannigfaltigkeit mit Rand}{}{}

\mathl{\partial M}{} und mit
\definitionsverweis {abzählbarer Basis der Topologie}{}{.}
Dann gibt es keine
\definitionsverweis {stetig differenzierbare Abbildung}{}{}
\maabbdisp {\varphi} { M } { \partial M
} {,}
deren
\definitionsverweis {Einschränkung}{}{}
auf $\partial M$ die Identität ist.}

 }

__NOEDITSECTION__

\inputaufgabeklausurloesung
{7}

{

Beweise den
\stichwort {Satz über die Lösung von gewöhnlichen Differentialgleichungen auf einer differenzierbaren Hyperfläche} {}

\mavergleichskette
{\vergleichskette
{ Y
}
{ \subseteq }{ \R^n
}
{  }{
}
{    }{
}
{   }{
}
}
{}{}{.}

}
{

Wir arbeiten mit der euklidischen Struktur auf dem $\R^n$ und mit dem
\definitionsverweis {Gradienten}{}{}
$\operatorname{Grad} \,  h$ zu $h$. Dieser ist nach Voraussetzung auf $Y$ nirgendwo gleich $0$ und dies überträgt sich wegen der Stetigkeit auf eine offene Umgebung von $Y$. Indem wir $U$ eventuell verkleinern, können wir annehmen, dass der Gradient auf ganz $U$ nullstellenfrei ist. Wir betrachten auf $U$  das neue Vektorfeld $G$, das durch

\mavergleichskettedisp
{\vergleichskette
{ G(P)
}
{ =} { F(P)-  { \frac{   \left\langle  F(P) ,
\operatorname{Grad} \,  h  (  P )    \right\rangle   }{  \Vert {
\operatorname{Grad} \,  h  (  P ) } \Vert^2   } }
\operatorname{Grad} \,  h  (  P )
}
{ } {
}
{ } {
}
{ } {
}
}
{}{}{}
gegeben ist. Für

\mavergleichskette
{\vergleichskette
{ P
}
{ \in }{ Y
}
{  }{
}
{    }{
}
{   }{
}
}
{}{}{}
ist

\mavergleichskette
{\vergleichskette
{ G(P)
}
{ = }{ F(P)
}
{  }{
}
{    }{
}
{   }{
}
}
{}{}{,}
da ja auf $Y$ der Gradient zu $h$ senkrecht auf dem Vektorfeld steht. Ferner besitzt $G$
\zusatzklammer {im Unterschied zu $F$} {} {}
die Eigenschaft, dass für alle Punkte

\mavergleichskette
{\vergleichskette
{ P
}
{ \in }{ U
}
{  }{
}
{    }{
}
{   }{
}
}
{}{}{}
der Vektor $G(P)$ senkrecht zum Gradienten steht, es ist ja

\mavergleichskettealignhandlinks
{\vergleichskettealignhandlinks
{ \left\langle  G(P) ,
\operatorname{Grad} \,  h  (  P )   \right\rangle
}
{ =} { \left\langle  F(P)- { \frac{   \left\langle  F(P) ,
\operatorname{Grad} \,  h  (  P )    \right\rangle  }{  \Vert {
\operatorname{Grad} \,  h  (  P ) } \Vert^2  } }
\operatorname{Grad} \,  h  (  P )  ,
\operatorname{Grad} \,  h  (  P )   \right\rangle
}
{ =} { \left\langle  F(P) ,
\operatorname{Grad} \,  h  (  P )   \right\rangle - \left\langle  F(P) ,
\operatorname{Grad} \,  h  (  P )   \right\rangle
}
{ =} { 0
}
{ } {
}
}
{}
{}{.}
Es sei nun
\maabbdisp {v} { I} { \R^n
} {}
die nach
[[Picard Lindelöf/Lokal Lipschitz/Lokale Existenz und Eindeutigkeit/Fakt|Satz 56.2 (Analysis (Osnabrück 2021-2023))]]
eindeutige Lösung zum Anfangswertproblem zu $G$ mit der Anfangsbedingung

\mavergleichskette
{\vergleichskette
{ v(t_0)
}
{ = }{ P
}
{ \in }{ Y
}
{    }{
}
{   }{
}
}
{}{}{.}
Dann ist

\mavergleichskettealign
{\vergleichskettealign
{ \left(Dh \circ v \right)_{t}
}
{ =} { \left(Dh  \right)_{v(t)} \circ  \left(D v \right)_{t}
}
{ =} { \left(Dh  \right)_{v(t)}  \left(  v'(t)  \right)
}
{ =} { \left(Dh  \right)_{v(t)}  \left(  G(v(t))  \right)
}
{ =} { \left\langle
\operatorname{Grad} \,  h  ( v(t) )  ,  G(v(t))    \right\rangle
}
}
{
\vergleichskettefortsetzungalign
{ =} { 0
}
{ } {}
{ } {}
{ } {}
}
{}{.}
Daher ist $h$ auf dem Bild $v( I)$ konstant und wegen

\mavergleichskette
{\vergleichskette
{ h(v(t_0))
}
{ = }{ c
}
{  }{
}
{    }{
}
{   }{
}
}
{}{}{}
ist

\mavergleichskette
{\vergleichskette
{ h(v(t))
}
{ = }{ c
}
{  }{
}
{    }{
}
{   }{
}
}
{}{}{}
für alle

\mavergleichskette
{\vergleichskette
{ t
}
{ \in }{ I
}
{  }{
}
{    }{
}
{   }{
}
}
{}{}{,}
also

\mavergleichskette
{\vergleichskette
{ v(t)
}
{ \in }{ Y
}
{  }{
}
{    }{
}
{   }{
}
}
{}{}{}
für alle

\mavergleichskette
{\vergleichskette
{ t
}
{ \in }{ I
}
{  }{
}
{    }{
}
{   }{
}
}
{}{}{.}

 }

__NOEDITSECTION__

\inputaufgabeklausurloesung
{4}

{

Es sei
\maabb {h} { I } { V
} {}
eine zweimal differenzierbare Kurve in einem
\definitionsverweis {euklidischen Vektorraum}{}{}
$V$. Zeige, dass bei

\mavergleichskette
{\vergleichskette
{h'(t)
}
{ \neq }{ 0
}
{  }{
}
{    }{
}
{   }{
}
}
{}{}{}
die Gleichheit

\mavergleichskettedisp
{\vergleichskette
{ \Vert { h'(t)} \Vert'
}
{ =} { { \frac{   \left\langle h'(t) , h^{\prime \prime} (t)  \right\rangle  }{  \left\langle h'(t) , h'(t)  \right\rangle  } } \Vert { h'(t)} \Vert
}
{ } {
}
{ } {
}
{ } {
}
}
{}{}{}
gilt.

}
{

Es sei

\mavergleichskette
{\vergleichskette
{u_1
}
{ = }{ h'(t)
}
{  }{
}
{    }{
}
{   }{
}
}
{}{}{,}
das wir zu einer Orthogonalbasis
\mathl{u_1,u_2 , \ldots , u_n}{} von $V$ ergänzen. Es seien
\mathl{h_1 , \ldots , h_n}{} die Koordinatenfunktionen von $h$ zu dieser Basis. Dann ist

\mavergleichskettealign
{\vergleichskettealign
{ \Vert { h'(t)} \Vert'
}
{ =} { \Vert { \begin{pmatrix} h_1'(t) \\\vdots\\ h_n'(t)   \end{pmatrix} } \Vert'
}
{ =} { \sqrt{  h_1'(t)^2 + \cdots + h_n'(t)^2 }'
}
{ =} { { \frac{   h_1(t) h_1'(t) + \cdots +  h_n(t) h_n'(t) }{  \sqrt{  h_1'(t)^2 + \cdots + h_n'(t)^2 }  } }
}
{ =} { { \frac{   h_1(t) h_1'(t) + \cdots +  h_n(t) h_n'(t) }{   h_1'(t)^2 + \cdots + h_n'(t)^2  } }  \sqrt{  h_1'(t)^2 + \cdots + h_n'(t)^2 }
}
}
{
\vergleichskettefortsetzungalign
{ =} { { \frac{   \left\langle h'(t) , h^{\prime \prime} (t)  \right\rangle   }{ \left\langle h'(t) , h'(t)  \right\rangle  } } \Vert { h'(t)} \Vert
}
{ } {
}
{ } {}
{ } {}
}
{}{.}

 }

__NOEDITSECTION__

\inputaufgabeklausurloesung
{0}

{

}
{  }

__NOEDITSECTION__

\inputaufgabeklausurloesung
{0}

{

}
{  }

__NOEDITSECTION__

\inputaufgabeklausurloesung
{0}

{

}
{  }

__NOEDITSECTION__

\inputaufgabeklausurloesung
{5}

{

Beweise die Aussage, dass die tangentiale Äquivalenz von Wegen auf einer Mannigfaltigkeit in einem Punkt $P$ mit einer beliebigen Karte überprüft werden kann.

}
{

Für eine
\definitionsverweis {differenzierbare Kurve}{}{}
\maabbdisp {\gamma} { I } { M
} {}
mit

\mavergleichskette
{\vergleichskette
{ 0
}
{ \in }{ I
}
{  }{
}
{    }{
}
{   }{
}
}
{}{}{}
und

\mavergleichskette
{\vergleichskette
{ \gamma(0)
}
{ = }{ P
}
{  }{
}
{    }{
}
{   }{
}
}
{}{}{}
und eine
\definitionsverweis {Karte}{}{}
\maabbdisp {\alpha} { U } { V
} {}
\zusatzklammer {mit

\mavergleichskettek
{\vergleichskettek
{P
}
{ \in }{ U
}
{  }{
}
{    }{
}
{   }{
}
}
{}{}{}
und

\mavergleichskette
{\vergleichskette
{V
}
{ \subseteq }{\R^n
}
{  }{
}
{    }{
}
{   }{
}
}
{}{}{}} {} {}
ändert sich der Ausdruck

\mathdisp {\left(  \alpha \circ \left(  \gamma \vert_{\gamma^{-1} (U)}  \right)  \right)'(0)} {  }

nicht, wenn man zu einem kleineren offenen Intervall

\mavergleichskette
{\vergleichskette
{ 0
}
{ \in }{ I'
}
{ \subseteq }{ I
}
{    }{
}
{   }{
}
}
{}{}{}
und einer kleineren offenen Menge

\mavergleichskette
{\vergleichskette
{P
}
{ \in }{ U'
}
{ \subseteq }{ U
}
{    }{
}
{   }{
}
}
{}{}{}
\zusatzklammer {mit der induzierten Karte} {} {}
übergeht. Wir können also davon ausgehen, dass
\mathkor {} {\gamma_1} {und} {\gamma_2} {}
auf dem gleichen Intervall definiert sind und ihre Bilder in $U$ liegen, und dass es für dieses $U$ zwei Karten
\maabbdisp {\alpha_1} { U } { V_1
} {}
und
\maabbdisp {\alpha_2} { U } { V_2
} {}
gibt. Dann folgt aus

\mavergleichskettedisp
{\vergleichskette
{ ( \alpha_1  \circ \gamma_1 )'(0)
}
{ =} { ( \alpha_1  \circ \gamma_2 )'(0)
}
{ } {
}
{ } {
}
{ } {
}
}
{}{}{}
nach
der [[Totale Differenzierbarkeit/K/Kettenregel/Fakt|Kettenregel]]
unter Verwendung der Differenzierbarkeit der
\definitionsverweis {Übergangsabbildung}{}{}

\mathl{\alpha_2 \circ \alpha_1^{-1}}{} sofort

\mavergleichskettealign
{\vergleichskettealign
{ \left(  \alpha_2 \circ \gamma_1  \right)'(0)
}
{ =} { \left(  \left(  \alpha_2 \circ \alpha_1^{-1}  \right) \circ \left(  \alpha_1 \circ \gamma_1  \right)  \right)'(0)
}
{ =} { \left(  D   \left(   \alpha_2 \circ \alpha_1^{-1}   \right)      \right)_{\alpha_1(P)} \left(   \left(  \alpha_1 \circ \gamma_1  \right)'(0)   \right)
}
{ =} { \left(  D   \left(   \alpha_2 \circ \alpha_1^{-1}   \right)      \right)_{\alpha_1(P)} \left(   \left(  \alpha_1 \circ \gamma_2  \right)'(0)   \right)
}
{ =} { \left(  \alpha_2 \circ \gamma_2  \right)'(0)
}
}
{}
{}{.}

 }

__NOEDITSECTION__

\inputaufgabeklausurloesung
{15 (3+3+4+5)}

{

Es sei

\mavergleichskette
{\vergleichskette
{T
}
{ = }{S^1 \times S^1
}
{  }{
}
{    }{
}
{   }{
}
}
{}{}{}
der Torus und

\mavergleichskette
{\vergleichskette
{S^2
}
{ \subset }{\R^3
}
{  }{
}
{    }{
}
{   }{
}
}
{}{}{}
die
\definitionsverweis {Einheitssphäre}{}{.}

a) Zeige, dass durch
\maabbeledisp {\varphi} {S^1 \times S^1} { S^2
} { \begin{pmatrix}  \alpha \\ \beta   \end{pmatrix} } { { \frac{   1 }{  2 } } \begin{pmatrix}  \sqrt{ 2-2 \sin \alpha   } \cos \beta  \\ \cos \alpha + \sin \beta \cos \alpha \\  1  + \sin \alpha -  \sin \beta + \sin \alpha \sin \beta   \end{pmatrix}
} {}
eine stetige Abbildung gegeben ist.

b) Zeige, dass $\varphi$ surjektiv ist.

c) Beschreibe die Fasern von $\varphi$.

d) Erläutere die Abbildung $\varphi$ unter Verwendung einer Skizze.

}
{

a) Wir zeigen zunächst, dass die Abbildung in der Tat auf der Sphäre landet, d.h. wir müssen zeigen, dass die Summe der Quadrate der drei Einträge gleich $1$ ist. Ohne Berücksichtigung des Vorfaktors ist dies

\mavergleichskettealign
{\vergleichskettealign
{
}
{ =} { 2  \cos^{ 2  }  \beta  -2 \sin \alpha  \cos^{ 2  }  \beta + \cos^{ 2  }  \alpha+ \cos^{ 2  }  \alpha \sin^{ 2  }  \beta +2 \cos^{ 2  }  \alpha\sin  \beta +1 + \sin^{ 2  } \alpha + \sin^{ 2  } \beta   + \sin^{ 2  } \alpha \sin^{ 2  } \beta +2 \sin \alpha  -2 \sin \beta   +2 \sin \alpha \sin \beta   -2 \sin \alpha \sin \beta +  2 \sin^{ 2  } \alpha \sin \beta  -2 \sin \alpha \sin^{ 2  } \beta
}
{ =} { 3 +   \cos^{ 2  }  \beta  -2 \sin \alpha  \cos^{ 2  }  \beta + \cos^{ 2  }  \alpha \sin^{ 2  }  \beta +2 \cos^{ 2  }  \alpha\sin  \beta  + \sin^{ 2  } \alpha \sin^{ 2  } \beta +2 \sin \alpha  -2 \sin \beta   +2 \sin^{ 2  } \alpha \sin \beta   - 2 \sin \alpha \sin^{ 2  } \beta
}
{ =} { 3 +   \cos^{ 2  }  \beta  +2 \sin \alpha   \left(  1 - \cos^{ 2  }  \beta   - \sin^{ 2  } \beta   \right)     + 2 \sin \beta \left(  - 1 + \cos^{ 2  }  \alpha    + \sin^{ 2  } \alpha  \right) + \cos^{ 2  }  \alpha\sin^{ 2  }  \beta  + \sin^{ 2  }  \alpha\sin^{ 2  }  \beta
}
{ =} { 3 +   \cos^{ 2  }  \beta  + \sin^{ 2  }  \beta
}
}
{
\vergleichskettefortsetzungalign
{ =} { 4
}
{ } {}
{ } {}
{ } {}
}
{}{.}

c) Die Berechnung der Faser über dem Nordpol
\mathl{\begin{pmatrix}  0 \\ 0\\  1   \end{pmatrix}}{} führt zu den Bedingungen

\mavergleichskettedisp
{\vergleichskette
{(1 - \sin \alpha ) \cos \beta
}
{ =} { 0
}
{ =} { \cos \alpha  (1 + \sin \beta )
}
{ } {
}
{ } {
}
}
{}{}{.}
Daher muss jeweils ein Faktor gleich $0$ sein. Bei

\mavergleichskettedisp
{\vergleichskette
{ \alpha
}
{ =} { { \frac{   1 }{  2 } } \pi
}
{ } {
}
{ } {
}
{ } {
}
}
{}{}{}
sind bei beliebigem $\beta$ beide Bedingungen erfüllt, ebenso sind bei

\mavergleichskettedisp
{\vergleichskette
{\beta
}
{ =} { { \frac{   3 }{  2 } } \pi
}
{ } {
}
{ } {
}
{ } {
}
}
{}{}{}
und beliebigem $\alpha$ beide Bedingungen erfüllt. Dabei ist wegen

\mavergleichskettedisp
{\vergleichskette
{ \sin  { \frac{   1 }{  2 } } \pi
}
{ =} { 1
}
{ } {
}
{ } {
}
{ } {
}
}
{}{}{}
bzw. wegen

\mavergleichskettedisp
{\vergleichskette
{ \sin  { \frac{   3 }{  2 } } \pi
}
{ =} { -1
}
{ } {
}
{ } {
}
{ } {
}
}
{}{}{}
die dritte Komponente gleich $1$. Bei

\mavergleichskette
{\vergleichskette
{\alpha
}
{ \neq }{ { \frac{   1 }{  2 } }  \pi
}
{  }{
}
{    }{
}
{   }{
}
}
{}{}{}
und

\mavergleichskette
{\vergleichskette
{\beta
}
{ \neq }{ { \frac{   3 }{  2 } } \pi
}
{  }{
}
{    }{
}
{   }{
}
}
{}{}{}
muss

\mavergleichskettedisp
{\vergleichskette
{ \cos \alpha
}
{ =} { 0
}
{ =} { \cos \beta
}
{ } {
}
{ } {
}
}
{}{}{}
sein. Die verbleibende Möglichkeit ist

\mavergleichskette
{\vergleichskette
{ \alpha
}
{ = }{ { \frac{   3 }{  2 } } \pi
}
{  }{
}
{    }{
}
{   }{
}
}
{}{}{}
und

\mavergleichskette
{\vergleichskette
{ \beta
}
{ = }{ { \frac{   1 }{  2 } } \pi
}
{  }{
}
{    }{
}
{   }{
}
}
{}{}{,}
doch dabei ist die dritte Komponente gleich $-1$.

d) Der Winkel $\alpha$ definiert den Punkt

\mavergleichskettedisp
{\vergleichskette
{P(\alpha)
}
{ =} { \begin{pmatrix}  0 \\ \cos \alpha \\  \sin \alpha   \end{pmatrix}
}
{ } {
}
{ } {
}
{ } {
}
}
{}{}{}
auf dem durch

\mavergleichskette
{\vergleichskette
{ x
}
{ = }{ 0
}
{  }{
}
{    }{
}
{   }{
}
}
{}{}{}
gegebenen Großreis. Der Nordpol und
\mathl{P(\alpha)}{} definieren den Halbierungspunkt

\mavergleichskettealign
{\vergleichskettealign
{Z(\alpha)
}
{ \defeq} { \begin{pmatrix}  0 \\ 0\\  1   \end{pmatrix} + { \frac{   1 }{  2 } }  \left(  P(\alpha)   - \begin{pmatrix}  0 \\ 0\\  1   \end{pmatrix}   \right)
}
{ =} { \begin{pmatrix}  0 \\ 0\\  1   \end{pmatrix} + { \frac{   1 }{  2 } }  \begin{pmatrix}  0 \\\cos \alpha \\ \sin \alpha - 1   \end{pmatrix}
}
{ =} { { \frac{   1 }{  2 } }  \begin{pmatrix}  0 \\\cos \alpha \\  1 + \sin \alpha   \end{pmatrix}
}
{ } {
}
}
{}
{}{.}
Der Bildpunkt
\mathl{P(\alpha, \beta)}{} wird auf dem Kreis $C$ mit Mittelpunkt
\mathl{Z(\alpha)}{} platziert, der durch den Nordpol und $P(\alpha)$ verläuft und der ganz auf der Sphäre liegt. Daher muss $C$ auf der Ebene liegen, die durch

\mavergleichskettedisp
{\vergleichskette
{P(\alpha) - \begin{pmatrix}  0 \\ 0\\  1   \end{pmatrix}
}
{ =} { { \frac{   1 }{  2 } } \begin{pmatrix}  0 \\\cos \alpha \\ \sin \alpha - 1   \end{pmatrix}
}
{ } {
}
{ } {
}
{ } {
}
}
{}{}{}
und
\mathl{\begin{pmatrix}  1 \\ 0\\  0   \end{pmatrix}}{} gegeben sind. Um mit der trigonometrischen Parametrisierung von $C$ arbeiten zu können, braucht man eine Orthonormalbasis, daher arbeiten wir mit

\mathdisp {O(\alpha) =  { \frac{   1 }{  \sqrt{ 2-2  \sin \alpha   }  } } \begin{pmatrix}  0 \\\cos \alpha \\ \sin \alpha - 1   \end{pmatrix}} { . }

Dies führt insgesamt auf

\mavergleichskettealign
{\vergleichskettealign
{\varphi(\alpha)
}
{ =} { Z (\alpha) + \cos  \beta \begin{pmatrix}  1 \\ 0\\  0   \end{pmatrix}  + \sin  \beta O(\alpha)
}
{ =} {  { \frac{   1 }{  2 } }  \begin{pmatrix}  0 \\\cos \alpha \\  1 + \sin \alpha   \end{pmatrix} + \cos  \beta { \frac{   \sqrt{ 2-2  \sin \alpha     }  }{  2  } } \begin{pmatrix}  1 \\ 0\\  0   \end{pmatrix}  + \sin  \beta { \frac{   \sqrt{ 2-2  \sin \alpha     }  }{  2  } } { \frac{   1 }{  \sqrt{ 2-2  \sin \alpha   }  } } \begin{pmatrix}  0 \\\cos \alpha \\ \sin \alpha - 1   \end{pmatrix}
}
{ =} { { \frac{   1 }{  2 } }       \begin{pmatrix}        \sqrt{ 2-2  \sin \alpha     }  \cos \beta         \\\cos \alpha \sin \beta   \\    1 + \sin \alpha  - \sin \beta  + \sin \alpha \sin \beta   \end{pmatrix}
}
{ } {
}
}
{}
{}{.}

 }

__NOEDITSECTION__

\inputaufgabeklausurloesung
{5}

{

Berechne das Wegintegral
\mathl{\int_\gamma \omega}{} zu
\maabbeledisp {\gamma} { [-1,0] } {\R^3
} { t } {(-t^2,t^3-1,t+2)
} {,}
für die $1$-Differentialform

\mavergleichskettedisp
{\vergleichskette
{ \omega
}
{ =} { x^3dx -yzdy +xz^2 dz
}
{ } {
}
{ } {
}
{ } {
}
}
{}{}{}
auf dem $\R^3$.

}
{

Es ist

\mavergleichskettealign
{\vergleichskettealign
{ \gamma^* \omega
}
{ =} { (-t^2)^3 d(-t^2) - (t^3-1) (t+2)d (t^3-1)  - t^2(t+2)^2d(t+2)
}
{ =} { 2 t^7 dt     - 3t^2 (t^4 +2t^3 -t -2) dt - t^2(  t^2  +4 t +4) dt
}
{ =} {  (2t^7 -3t^6 -6t^5   +3t^3 +6t^2 -t^4  -4t^3 -4t^2       )dt
}
{ =} { (2t^7 -3t^6 -6t^5   -t^4  -t^3 +2 t^2       )dt
}
}
{}
{}{.}
Daher ist

\mavergleichskettealign
{\vergleichskettealign
{ \int_\gamma \omega
}
{ =} { \int_{-1}^0    (2t^7 -3t^6 -6t^5   -t^4  -t^3 +2 t^2       )dt
}
{ =} { ({ \frac{   1 }{  4 } }  t^8 - { \frac{   3 }{  7 } }  t^7 - t^6- { \frac{   1 }{  5 } } t^5 - { \frac{   1 }{  4 } } t^4 + { \frac{   2 }{  3 } } t^3) \vert_{  -1 }^{  0 }
}
{ =} { -( { \frac{   1 }{  4 } }  + { \frac{   3 }{  7 } }  -1 + { \frac{   1 }{  5 } }  - { \frac{   1 }{  4 } }  - { \frac{   2 }{  3 } } )
}
{ =} { - { \frac{   3 }{  7 } }  +1 - { \frac{   1 }{  5 } } + { \frac{   2 }{  3 } }
}
}
{
\vergleichskettefortsetzungalign
{ =} { { \frac{   -45 +105-21+70 }{  105 } }
}
{ =} { { \frac{   109 }{  105 } }
}
{ } {}
{ } {}
}
{}{.}

 }

__NOEDITSECTION__

\inputaufgabeklausurloesung
{0}

{

}
{  }

__NOEDITSECTION__

\inputaufgabeklausurloesung
{3 (2+1)}

{

Es sei
\maabbeledisp {f} {\R} {\R
} { t } {f(t)
} {,}
die durch

\mavergleichskettedisp
{\vergleichskette
{ f(t)
}
{ =} { { \frac{   \sin^{ 3  } (t^4)  }{  1+t^2 } }
}
{ } {
}
{ } {
}
{ } {
}
}
{}{}{}
gegeben ist.

a) Berechne die äußere Ableitung von $f$.

b) Berechne die äußere Ableitung von $fdt$.

}
{

a) Es ist
\mathl{df=f'dt}{,} und zwar ist nach der Quotientenregel

\mavergleichskettealign
{\vergleichskettealign
{f'
}
{ =} { { \frac{   (1+t^2)  3\sin^{ 2  } (t^4) \cos (t^4) 4 t^3  - 2t \sin^{ 3  } (t^4)  }{ (1+t^2)^2 } }
}
{ =} { { \frac{   12 (t^3+t^5) \sin^{ 2  } (t^4) \cos (t^4)   - 2t \sin^{ 3  } (t^4)  }{  1+2t^2+t^4 } }
}
{ } {
}
{ } {
}
}
{}
{}{.}

b) Die äußere Ableitung von $fdt$ ist
\mathl{f'dt \wedge dt =0}{.}

 }

__NOEDITSECTION__

\inputaufgabeklausurloesung
{3}

{

Man gebe ein Beispiel für eine stetig differenzierbare Abbildung
\maabbdisp {\varphi} { H } { H
} {}
einer euklidischen Halbebene in sich mit der Eigenschaft, dass genau ein Randpunkt von $H$ in einen Randpunkt und alle anderen Punkte in das Innere der Halbebene abgebildet werden.

}
{

Wir betrachten die positive Halbebene

\mavergleichskettedisp
{\vergleichskette
{ H
}
{ =} { \left\{ (x,y)  \mid   x \geq 0        \right\}
}
{ } {
}
{ } {
}
{ } {
}
}
{}{}{}
und die Abbildung
\maabbeledisp {\varphi} { H } {H
} {(x,y)} { \left(  x+y^2  , \,   y                   \right)
} {.}
Als polynomiale Abbildung ist $\varphi$ stetig differenzierbar. Der Randpunkt
\mathl{(0,0)}{} wird auf den Randpunkt
\mathl{(0,0)}{} abgebildet. Bei allen anderen Punkten von $H$ ist

\mavergleichskette
{\vergleichskette
{ x
}
{ > }{  0
}
{  }{
}
{    }{
}
{   }{
}
}
{}{}{}
oder

\mavergleichskette
{\vergleichskette
{ y
}
{ \neq }{ 0
}
{  }{
}
{    }{
}
{   }{
}
}
{}{}{}
und daher stets

\mavergleichskettedisp
{\vergleichskette
{ x+y^2
}
{ >} {  0
}
{ } {
}
{ } {
}
{ } {
}
}
{}{}{,}
die erste Komponente ist also positiv.

 }

__NOEDITSECTION__

\inputaufgabeklausurloesung
{0}

{

}
{  }

__NOEDITSECTION__

\inputaufgabeklausurloesung
{10}

{

Beweise die Quaderversion des Satzes von Stokes.

}
{

Da beide Seiten dieser Gleichung linear in $\omega$ sind, können wir annehmen, dass $\omega$ die Gestalt

\mavergleichskettedisp
{\vergleichskette
{ \omega
}
{ =} { f dx_2 \wedge \ldots \wedge dx_n
}
{ } {
}
{ } {
}
{ } {
}
}
{}{}{}
mit einer in einer offenen Umgebung von $Q$ definierten stetig differenzierbaren Funktion $f$ besitzt. Die Integrale sind links und rechts
\definitionsverweis {Lebesgue-Integrale}{}{}
zu stetigen Funktionen auf Teilmengen des
\mathkor {} {\R^n} {bzw.} {\R^{n-1}} {.}
Daher können wir auf beiden Seiten zum
\definitionsverweis {topologischen Abschluss}{}{}
übergehen, da dadurch die in Frage stehenden Integrationsbereiche nur um eine Nullmenge verändert werden, sodass dies die Integrale nicht ändert.

Wir schreiben den abgeschlossenen Quader als

\mavergleichskettedisp
{\vergleichskette
{ \overline{Q}
}
{ =} { [a,b] \times \tilde{Q}
}
{ } {
}
{ } {
}
{ } {
}
}
{}{}{.}
Wir wenden
[[Differentialform/Lokal/Zurückziehen unter partiell konstanter Abbildung/Fakt|Korollar 14.10 (Differentialgeometrie (Osnabrück 2023))]]
auf jede Seite $S$ ausgenommen
\mathkor {} {a \times \tilde{Q}} {und} {b \times \tilde{Q}} {}
an und erhalten darauf

\mavergleichskettedisp
{\vergleichskette
{
\int_{ S  }   \omega
}
{ =} {
\int_{  S  }  f dx_2 \wedge \ldots \wedge dx_n
}
{ =} {
\int_{ S  }   0
}
{ =} { 0
}
{ } {
}
}
{}{}{,}
da auf diesen Seiten jeweils eine der Variablen
\mathl{x_2 , \ldots , x_n}{} konstant ist. Aufgrund
des [[Sigmaendliche Räume/Satz von Fubini/Fakt|Satzes von Fubini]]
und
des [[Riemann-Integral/Hauptsatz/Newton-Leibniz/Fakt|Hauptsatzes der Infinitesimalrechnung]]
\zusatzklammer {angewendet auf jedes fixierte \mathlk{(x_2 , \ldots , x_n)}{}} {} {}
gilt

\mavergleichskettealignhandlinks
{\vergleichskettealignhandlinks
{
\int_{ Q  }   d\omega
}
{ =} {
\int_{  Q  }   df \wedge  dx_2 \wedge \ldots \wedge dx_n
}
{ =} {
\int_{  Q  }   \left(  \sum_{j = 1}^n { \frac{   \partial   f   }{  \partial   x_j   } } dx_j  \right) \wedge  dx_2 \wedge \ldots \wedge dx_n
}
{ =} {
\int_{  Q  }   { \frac{   \partial   f   }{  \partial   x_1   } } dx_1 \wedge  dx_2 \wedge \ldots \wedge dx_n
}
{ =} {
\int_{ \tilde{Q}   }   \left(  \int_{[a,b] } { \frac{   \partial   f   }{  \partial   x_1   } } dx_1  \right)  dx_2 \wedge \ldots \wedge dx_n
}
}
{
\vergleichskettefortsetzungalign
{ =} {
\int_{  \tilde{Q}   }   \left(  f(b,x_2 , \ldots , x_n ) - f(a,x_2 , \ldots , x_n )   \right) dx_2 \wedge \ldots \wedge dx_n
}
{ =} {
\int_{ b \times \tilde{Q}   }   f dx_2 \wedge \ldots \wedge dx_n  -
\int_{  a \times \tilde{Q}   }  f dx_2 \wedge \ldots \wedge dx_n
}
{ =} { \sum_{S \text{ orientierte Seite von } Q}
\int_{  S  }  f dx_2 \wedge \ldots \wedge dx_n
}
{ =} {
\int_{  \partial Q  }   f dx_2 \wedge \ldots \wedge dx_n
}
}
{
\vergleichskettefortsetzungalign
{ =} {
\int_{  \partial Q  }   \omega
}
{ } {}
{ } {}
{ } {}
}{.}

 }

__NOEDITSECTION__

\inputaufgabeklausurloesung
{4}

{

Bestimme für die Kurve
\maabbeledisp {\psi} {\R} { {\mathbb H}
} { t } { \left(  t   , \,  e^t                   \right)
} {,}
in die
\definitionsverweis {Halbebene von Poincaré}{}{}
${\mathbb H}$ die
\definitionsverweis {Christoffelsymbole}{}{}
für den
\definitionsverweis {zurückgezogenen Zusammenhang}{}{}
\zusatzklammer {zum
\definitionsverweis {Levi-Civita-Zusammenhang}{}{}
auf ${\mathbb H}$} {} {}
auf $\R$.

}
{

Die einzige Standardableitung auf $\R$  ist $\partial$, deshalb verzichten wir auf einen unteren Index für die Christoffelsymbole. Nach
[[Linearer Zusammenhang/Vertikale Ableitung/Längs Abbildung/Eigenschaften/Fakt|Fakt]] [[Kurs:Differentialgeometrie/Linearer Zusammenhang/Vertikale Ableitung/Längs Abbildung/Eigenschaften/Fakt/Faktreferenznummer|*****]]  (2)
ist für

\mavergleichskette
{\vergleichskette
{ \ell,k
}
{ = }{ 1,2
}
{  }{
}
{    }{
}
{   }{
}
}
{}{}{}

\mavergleichskettedisp
{\vergleichskette
{ \tilde{\Gamma}_\ell^k (t)
}
{ =} { \psi_1'(t) \Gamma_{1 \ell}^k (\psi(t)) + \psi'_2(t) \Gamma_{2 \ell}^k (\psi(t))
}
{ =} { \Gamma_{1 \ell}^k (t,e^t)   + e^t \Gamma_{2 \ell}^k (t,e^t)
}
{ } {}
{ } {}
}
{}{}{.}
Mit
[[Halbebene/Poincaré/Riemannsche Geometrie/Christoffelsymbole/Beispiel|Beispiel]] [[Kurs:Differentialgeometrie/Halbebene/Poincaré/Riemannsche Geometrie/Christoffelsymbole/Beispiel/Beispielreferenznummer|*****]]
ergibt sich

\mavergleichskettedisp
{\vergleichskette
{ \tilde{\Gamma}_1^1 (t)
}
{ =} { \Gamma_{1 1}^1 (t,e^t) + e^t \Gamma_{2 1}^1 (t,e^t)
}
{ =} {  -e^t e^{-t}
}
{ =} { -1
}
{ } {
}
}
{}{}{,}

\mavergleichskettedisp
{\vergleichskette
{ \tilde{\Gamma}_1^2 (t)
}
{ =} { \Gamma_{1 1}^2 (t,e^t) + e^t \Gamma_{2 1}^2 (t,e^t)
}
{ =} { e^{-t}
}
{ } {
}
{ } {
}
}
{}{}{,}

\mavergleichskettedisp
{\vergleichskette
{ \tilde{\Gamma}_2^1 (t)
}
{ =} { \Gamma_{1 2}^1 (t,e^t) + e^t \Gamma_{2 2}^1 (t,e^t)
}
{ =} {  -e^{-t}
}
{ } {
}
{ } {
}
}
{}{}{}
und

\mavergleichskettedisp
{\vergleichskette
{ \tilde{\Gamma}_2^2 (t)
}
{ =} { \Gamma_{1 2}^2 (t,e^t) + e^t \Gamma_{2 2}^2 (t,e^t)
}
{ =} { - e^t e^{-t}
}
{ =} { -1
}
{ } {
}
}
{}{}{.}

 }

 [[Kategorie:Latexseite]]
